%==============================================================
% BAB I  PENDAHULUAN
%==============================================================

\chapter{PENDAHULUAN}

%--------------------------------------------------------------
\section{Latar Belakang}
%--------------------------------------------------------------

\textit{Retinopathy of Prematurity} (ROP) adalah gangguan pembuluh darah retina yang terjadi
pada bayi yang lahir sebelum waktunya. Kondisi ini muncul akibat pertumbuhan pembuluh darah
retina yang tidak sempurna setelah kelahiran prematur, dan dapat berkembang menjadi ablasio
retina yang menyebabkan kebutaan permanen jika tidak ditangani sejak awal
\citep{Chiang2021, Vahidmoghadam2026}. Kemajuan di bidang perawatan neonatal telah meningkatkan
angka kelangsungan hidup bayi prematur di seluruh dunia, sehingga jumlah kasus ROP yang
memerlukan pemantauan juga meningkat secara signifikan.

Secara klinis, keparahan ROP dibagi ke dalam beberapa \textit{stage} berdasarkan perubahan
struktural yang terlihat pada retina melalui pemeriksaan fundus. \textit{Stage} 1 ditandai
dengan adanya garis batas (\textit{demarcation line}) yang memisahkan area retina yang sudah
dan belum memiliki pembuluh darah. \textit{Stage} 2 menunjukkan penebalan garis tersebut
menjadi tonjolan jaringan (\textit{ridge}). Pada \textit{stage} 3, terjadi pertumbuhan
pembuluh darah baru ke arah vitreous yang disebut neovaskularisasi ekstraretina
\citep{Chiang2021}. Diagnosis yang tepat pada \textit{stage} awal sangat penting karena
intervensi medis yang tepat waktu dapat mencegah perkembangan ke kondisi yang lebih parah.

Diagnosis ROP saat ini bergantung pada pemeriksaan langsung oleh dokter spesialis mata
menggunakan kamera fundus khusus. Tantangan utama dalam proses ini adalah kelangkaan dokter
spesialis retina anak, terutama di negara berkembang, serta tingginya variabilitas antar-dokter
dalam menilai \textit{stage} penyakit \citep{Tong2020}. Otomatisasi melalui \textit{computer
vision} menawarkan potensi untuk membantu proses skrining massal dan mengurangi beban tenaga
medis.

Namun, klasifikasi otomatis \textit{stage} ROP menghadapi tantangan tersendiri. Citra fundus
bayi memiliki kualitas yang lebih rendah dibanding citra fundus dewasa karena adanya reflek
kornea, pupil yang lebih kecil, dan kesulitan dalam mempertahankan fiksasi pandangan bayi
selama pengambilan gambar. Selain itu, perbedaan visual antar-\textit{stage}, khususnya antara
\textit{Stage} 1 dan \textit{Stage} 2, sangat halus sehingga bahkan model jaringan saraf yang
telah dilatih dengan data besar pun sering kesulitan membedakan keduanya \citep{Zhao2024}.

Penelitian sebelumnya menunjukkan bahwa teknik \textit{enhancement} citra seperti
\textit{contrast-limited adaptive histogram equalization} (CLAHE) dan koreksi pencahayaan dapat
membantu menonjolkan struktur klinis penting pada citra fundus bayi \citep{Rahim2024}.
Demikian pula, segmentasi pembuluh darah secara otomatis telah terbukti dapat menghasilkan
peta spasial yang berguna sebagai panduan bagi model klasifikasi \citep{Almeida2024}. Namun,
pengaruh masing-masing pendekatan terhadap performa klasifikasi \textit{stage} ROP secara
spesifik belum banyak diteliti secara sistematis dan terkontrol.

Penelitian ini dirancang untuk menjawab pertanyaan tersebut secara langsung melalui sebuah
\textit{pipeline} yang membandingkan berbagai bentuk representasi citra---mulai dari piksel
RGB mentah, citra hasil \textit{enhancement}, hingga peta pembuluh darah---sebagai masukan
untuk model jaringan saraf yang sama. Penelitian ini juga menyelidiki apakah sebuah jaringan
saraf kecil yang dilatih dari awal (\textit{from scratch}) tanpa bobot pra-latih dapat
mencapai performa yang bersaing dengan pendekatan berbasis fitur klasik.

%--------------------------------------------------------------
\section{Rumusan Masalah}
%--------------------------------------------------------------

Penelitian ini berfokus pada tiga pertanyaan teknis utama:

\vspace{2mm}
\noindent\textbf{Apakah \textit{enhancement} citra meningkatkan akurasi klasifikasi?}
Koreksi pencahayaan dan CLAHE pada ruang warna L*a*b* diuji untuk melihat apakah
perbaikan kontras citra secara visual berkorelasi dengan peningkatan performa model
pada tugas klasifikasi \textit{stage} ROP.

\vspace{2mm}
\noindent\textbf{Apakah struktur pembuluh darah saja sudah cukup untuk membedakan \textit{stage}?}
\textit{Vessel mask} biner hasil segmentasi pembuluh darah diuji secara terpisah untuk
mengukur sejauh mana informasi bentuk dan distribusi pembuluh darah dapat membantu membedakan
antara \textit{stage} 1 hingga 3 tanpa bergantung pada tekstur retina lainnya.

\vspace{2mm}
\noindent\textbf{Apakah representasi \textit{softmap} spasial lebih baik dari ringkasan fitur klasik?}
Model jaringan saraf yang menerima peta spasial (\textit{softmap}) pembuluh darah dan area
tepi secara langsung dibandingkan dengan model berbasis fitur skalar klasik yang menggunakan
informasi vaskular yang sama, untuk menguji apakah pemrosesan spasial oleh jaringan saraf
memberikan keuntungan yang nyata.

%--------------------------------------------------------------
\section{Tujuan}
%--------------------------------------------------------------

Tujuan penelitian ini adalah:

\begin{enumerate}
    \item Membangun dan mengevaluasi \textit{pipeline} segmentasi pembuluh darah retina secara
    klasik pada citra fundus bayi menggunakan kombinasi filter Gabor dan Meijering.
    \item Membandingkan performa model \textit{TinyResNetV2} dalam mengklasifikasi
    \textit{stage} ROP menggunakan empat skenario representasi input: citra RGB mentah,
    citra hasil \textit{enhancement}, \textit{vessel mask} biner, dan citra \textit{enhanced}
    yang dipandu \textit{vessel mask}.
    \item Menguji apakah representasi \textit{softmap} spasial sebagai masukan model jaringan
    saraf dapat melampaui pendekatan berbasis fitur klasik yang menggunakan informasi
    vaskular yang sama.
    \item Mengevaluasi dampak protokol validasi silang yang mempertimbangkan kelompok citra
    (\textit{group-aware cross-validation}) terhadap keandalan estimasi performa model.
\end{enumerate}

%--------------------------------------------------------------
\section{Ruang Lingkup}
%--------------------------------------------------------------

Batasan dan lingkup penelitian ini ditetapkan sebagai berikut:

\begin{itemize}
    \item \textbf{Dataset Utama:} Menggunakan dataset publik Zhao2024 \citep{Zhao2024} yang
    terdiri dari 1.099 citra fundus bayi dalam lima kelas untuk eksperimen klasifikasi utama.
    \item \textbf{Kelas Klasifikasi:} Mencakup Normal, Stage 1, Stage 2, Stage 3, dan
    \textit{laser scars}. Kelas \textit{laser scars} dimasukkan dalam klasifikasi akhir
    untuk mendekati pengaturan 5-kelas yang digunakan dalam publikasi Zhao2024.
    \item \textbf{Dataset Sekunder:} Menggunakan bagian HVDROPDB-BV dari Agrawal2021
    \citep{Agrawal2021} khusus untuk mengevaluasi kualitas segmentasi pembuluh darah secara
    kuantitatif menggunakan metrik Dice.
    \item \textbf{Model:} Menggunakan arsitektur \textit{custom TinyResNetV2} yang dilatih
    dari awal tanpa bobot pra-latih. Tidak ada penggunaan model pra-latih pada \textit{dataset}
    seperti ImageNet dalam penelitian ini.
    \item \textbf{Konteks:} Penelitian ini bersifat akademis dan bukan merupakan sistem
    diagnostik klinis. Hasil yang dilaporkan tidak boleh digunakan untuk pengambilan keputusan
    medis.
\end{itemize}

%--------------------------------------------------------------
\section{Manfaat Penelitian}
%--------------------------------------------------------------

Penelitian ini memberikan beberapa kontribusi yang dapat digunakan oleh peneliti dan
pengembang sistem di bidang analisis citra medis:

\begin{itemize}
    \item Menunjukkan secara eksperimental bahwa representasi spasial dari informasi pembuluh
    darah lebih efektif daripada ringkasan fitur skalar dalam tugas klasifikasi \textit{stage}
    ROP.
    \item Memberikan bukti empiris tentang pentingnya protokol validasi silang yang
    mempertimbangkan kelompok citra serupa (\textit{group-aware}) agar estimasi performa
    yang dilaporkan mencerminkan kemampuan generalisasi model secara realistis.
    \item Mendokumentasikan pengembangan sistematis \textit{pipeline} segmentasi pembuluh
    darah klasik melalui serangkaian konfigurasi yang dievaluasi secara terkontrol.
    \item Menyediakan kerangka percobaan yang dapat direplikasi untuk penelitian lanjutan
    tentang klasifikasi \textit{stage} ROP dengan data yang lebih besar atau arsitektur yang
    lebih canggih.
\end{itemize}
